\documentclass[12pt,a4paper]{article}

% ── Encodage et langue ──────────────────────────────────────────
\usepackage[utf8]{inputenc}
\usepackage[T1]{fontenc}
\usepackage[french]{babel}

% ── Mise en page ─────────────────────────────────────────────────
\usepackage[top=2.5cm, bottom=2.5cm, left=2.5cm, right=2.5cm]{geometry}
\usepackage{setspace}
\onehalfspacing

% ── Mathématiques ────────────────────────────────────────────────
\usepackage{amsmath, amssymb, amsthm}
\usepackage{mathtools}
\usepackage{bm}

% ── Tableaux ─────────────────────────────────────────────────────
\usepackage{booktabs}
\usepackage{tabularx}
\usepackage{array}
\usepackage{longtable}
\usepackage{multirow}
\usepackage{rotating}
\usepackage{pdflscape}

% ── Typographie ──────────────────────────────────────────────────
\usepackage{microtype}
\usepackage{lmodern}
\usepackage{parskip}
\setlength{\parindent}{0pt}
\setlength{\parskip}{6pt}

% ── Liens ────────────────────────────────────────────────────────
\usepackage[colorlinks=true, linkcolor=black, citecolor=black,
            urlcolor=black]{hyperref}

% ── En-tête / pied de page ───────────────────────────────────────
\usepackage{fancyhdr}
\pagestyle{fancy}
\fancyhf{}
\fancyhead[L]{\small\textit{BiVAT  État de l'art \& Justification architecturale}}
\fancyhead[R]{\small\textit{MOUKOKO EKAMBI G.M.}}
\fancyfoot[C]{\thepage}
\renewcommand{\headrulewidth}{0.4pt}

% ── Sections ─────────────────────────────────────────────────────
\usepackage{titlesec}
\titleformat{\section}{\large\bfseries}{\thesection.}{0.5em}{}
\titleformat{\subsection}{\normalsize\bfseries}{\thesubsection}{0.5em}{}
\titleformat{\subsubsection}{\normalsize\itshape}{\thesubsubsection}{0.5em}{}

% ── Commandes utilitaires ────────────────────────────────────────
\newcommand{\nd}{--}

% ── Environnement encart neutre ──────────────────────────────────
\usepackage{mdframed}
\newmdenv[
  linewidth=0.4pt,
  linecolor=black,
  backgroundcolor=white,
  innerleftmargin=6pt,
  innerrightmargin=6pt,
  innertopmargin=4pt,
  innerbottommargin=4pt,
  skipabove=6pt,
  skipbelow=6pt
]{encart}

% ── Bibliographie ────────────────────────────────────────────────
\usepackage{natbib}

% ════════════════════════════════════════════════════════════════
\begin{document}

% ── Page de titre ────────────────────────────────────────────────
\begin{titlepage}
\centering
\vspace*{2cm}

{\LARGE\bfseries
DÉTECTION DE VARIATIONS ATYPIQUES DANS LES SÉRIES TEMPORELLES
FINANCIÈRES PAR ARCHITECTURES HYBRIDES VAE-TRANSFORMER \\[0.4em]
\large\normalfont État de l'art outillé et justification architecturale
de BiVAT\\[0.4em]
}

\vspace{2cm}

{\large\bfseries MOUKOKO EKAMBI GEORGES MIGUEL}\\[0.5em]
{\normalsize\itshape
Stagiaire, Banque des États de l'Afrique Centrale (BEAC)\\
M2 Data Science \& IA  ENSPD / Université de Douala, GISDIA Laboratory
}

\vspace{2cm}

{\normalsize Août 2026}

\vfill

{\small
Document rédigé dans le cadre du stage de recherche à la Banque des États
de l'Afrique Centrale (BEAC), M2 Data Science \& IA,
ENSPD / Université de Douala  GISDIA Laboratory.
}
\end{titlepage}

% ── Résumé ───────────────────────────────────────────────────────
\begin{abstract}
La détection non supervisée de variations atypiques dans des séries temporelles
financières multivariées constitue un problème fondamental pour la surveillance
macroprudentielle. Ce document présente un état de l'art exhaustif et outillé
des architectures d'intelligence artificielle mobilisées à cette fin, enrichi
de benchmarks normalisés sur cinq datasets de référence (SMD, MSL, SMAP,
SWaT, PSM) et d'une matrice qualitative comparative sur dix modèles. Sur la
base d'une cartographie de trois lacunes structurelles non résolues
simultanément dans la littérature, quatre choix architecturaux sont justifiés
empiriquement pour la méthode proposée BiVAT (\textit{Bidirectional
VAE-Associative Transformer}) : un encodeur hybride Transformer bidirectionnel
et BiLSTM, la Conformal Prediction pour la quantification d'incertitude, SHAP
pour l'explicabilité réglementaire, et un critère discriminant hybride combinant
reconstruction probabiliste VAE et Association Discrepancy. Ce dernier apporte
un gain documenté de $+18{,}82$ points de F1 moyen par rapport à la
reconstruction seule (ablation study, Xu et al.\ 2022). Le dispositif est
dimensionné pour le dataset BEAC-CEMAC (195 observations mensuelles, 55
indicateurs IFS) et satisfait aux exigences d'auditabilité algorithmique COBAC.

\bigskip
\noindent\textbf{Mots-clés :} détection d'anomalies, séries temporelles
multivariées, autoencodeur variationnel, Transformer, BiLSTM, Association
Discrepancy, Conformal Prediction, SHAP, BEAC, surveillance macroprudentielle.
\end{abstract}

\tableofcontents
\newpage

% ════════════════════════════════════════════════════════════════
\section{Introduction}

\subsection{Contexte général}

La détection automatique de variations atypiques désigne le processus
d'identification de patrons dans des données qui s'écartent statistiquement du
comportement attendu. Dans le domaine monétaire et financier, ces variations
peuvent signaler des erreurs de collecte, des chocs macroéconomiques exogènes,
des ruptures structurelles dans les dynamiques de crédit ou de dépôt, ou des
comportements frauduleux. Pour une institution comme la Banque des États de
l'Afrique Centrale (BEAC), dont la mission est d'assurer la stabilité monétaire
dans la zone CEMAC, la détection précoce et fiable de ces variations constitue
une exigence opérationnelle et réglementaire de premier ordre. La conception de
tout système de détection automatique destiné à la BEAC est soumise à une
contrainte supplémentaire non négligeable : l'exigence d'auditabilité
algorithmique imposée par le cadre de surveillance prudentielle COBAC, qui
requiert que toute décision automatisée soit traçable, explicable et
reproductible.

Le dataset de référence de ce travail  les séries monétaires des Autres
Sociétés de Dépôts (2SR) selon la nomenclature IFS du FMI  présente des
caractéristiques structurelles qui en font un terrain d'expérimentation
exigeant : environ 195 observations mensuelles (décembre 2001 à mars 2026),
environ 55 indicateurs financiers multivariés couvrant l'actif et le passif
du bilan des institutions de dépôt du Cameroun, une forte hétérogénéité
d'échelles (en milliards de FCFA), des valeurs manquantes sporadiques, une
saisonnalité budgétaire marquée et des ruptures structurelles documentées
(choc Covid-19 en 2020--2021, dynamique de crédit 2024).

\subsection{Pourquoi le cadre non supervisé ?}

En l'absence d'étiquettes d'anomalies historiquement certifiées par un processus
d'audit formalisé, les approches supervisées sont inapplicables directement. Le
cadre non supervisé, fondé sur la modélisation de la distribution normale et la
détection par déviation, est le seul praticable. Parmi les familles de méthodes
disponibles, les architectures profondes à base d'autoencodeurs variationnels
(VAE) et de Transformers présentent la cohérence structurelle la plus forte avec
les caractéristiques du dataset BEAC pour trois raisons : leur capacité à
modéliser des distributions multivariées complexes sans hypothèses paramétriques
fortes ; leur aptitude à capturer les dépendances temporelles à longue portée ;
et leur extensibilité vers des scores probabilistes d'anomalie intégrant
l'incertitude du modèle.

\subsection{Organisation de ce document}

Ce document est organisé en onze sections. Les sections 2 à 6 constituent
l'état de l'art technique, organisé en quatre axes (fondations théoriques,
Transformers, hybrides VAE-Transformer, architectures avancées), enrichi pour
chaque modèle d'un encart de positionnement benchmark et d'une grille qualitative.
La section~7 présente la synthèse transversale avec le tableau de benchmarks
normalisés, la matrice qualitative comparative et la cartographie des lacunes
ancrée empiriquement. La section~8 constitue la contribution principale de ce
document : la justification multi-critères des quatre choix architecturaux de
BiVAT, appuyée sur des ablation studies publiées, des comparaisons empiriques
et des références réglementaires institutionnelles. Les sections~9 à 11
présentent respectivement l'architecture BiVAT, sa cohérence avec le dataset
BEAC, et la conclusion.

% ════════════════════════════════════════════════════════════════
\section{Cadre Théorique Fondateur}

\subsection{Inférence variationnelle et ELBO}

Soit un ensemble d'observations $\mathcal{X} = \{\mathbf{x}^{(i)}\}_{i=1}^{N}$
généré par un processus impliquant une variable latente non observée
$\mathbf{z} \in \mathbb{R}^d$. Le modèle génératif spécifie :

$$p_\theta(\mathbf{x}, \mathbf{z}) = p_\theta(\mathbf{x}|\mathbf{z})\,p(\mathbf{z})$$

La vraisemblance marginale
$p_\theta(\mathbf{x}) = \int p_\theta(\mathbf{x}|\mathbf{z})\,p(\mathbf{z})\,
d\mathbf{z}$ est intractable dès que $p_\theta(\mathbf{x}|\mathbf{z})$ est un
réseau de neurones. L'inférence variationnelle contourne cette intractabilité
en introduisant une distribution approchée $q_\phi(\mathbf{z}|\mathbf{x})$ et
en maximisant l'Evidence Lower BOund (ELBO) :

$$\log p_\theta(\mathbf{x}) \;\geq\; \mathcal{L}(\theta, \phi;\,\mathbf{x})
= \underbrace{\mathbb{E}_{q_\phi(\mathbf{z}|\mathbf{x})}\!\left[
\log p_\theta(\mathbf{x}|\mathbf{z})\right]}_{\text{terme de reconstruction}}
- \underbrace{D_{\mathrm{KL}}\!\left(q_\phi(\mathbf{z}|\mathbf{x})
\;\|\; p(\mathbf{z})\right)}_{\text{terme de régularisation KL}}$$

\subsection{Le mécanisme d'attention multi-têtes}

Vaswani et al.\ (2017) proposent l'opération d'auto-attention à produit scalaire :

$$\mathrm{Attention}(\mathbf{Q}, \mathbf{K}, \mathbf{V}) =
\mathrm{Softmax}\!\left(\frac{\mathbf{Q}\mathbf{K}^\top}{\sqrt{d_k}}\right)\mathbf{V}$$

L'attention multi-têtes concatène $H$ têtes parallèles :

$$\mathrm{MHA}(\mathbf{Q},\mathbf{K},\mathbf{V}) =
\mathrm{Concat}(\mathrm{head}_1, \ldots, \mathrm{head}_H)\,\mathbf{W}^O$$

La matrice d'attention $\mathbf{A} \in \mathbb{R}^{T\times T}$ encode, pour
chaque instant $t$, une distribution de poids sur l'ensemble des instants de
la séquence. Dans le cadre non causal (bidirectionnel), aucun masquage n'est
appliqué : chaque instant peut s'associer librement à tout autre instant passé
ou futur. C'est cette propriété qui fonde le choix d'un encodeur Transformer
non causal dans BiVAT.

% ════════════════════════════════════════════════════════════════
\section{Axe 1 : Fondations Théoriques}

\subsection{[F1] Kingma \& Welling (2013)  Auto-Encoding Variational Bayes}

\textbf{Problème formel.} Apprendre conjointement une représentation latente
compacte et probabiliste $q_\phi(\mathbf{z}|\mathbf{x})$ et un modèle génératif
$p_\theta(\mathbf{x}|\mathbf{z})$ de manière non supervisée, en rendant le
processus d'optimisation différentiable end-to-end malgré l'opération
d'échantillonnage.

\textbf{Architecture et paramétrisation.} L'encodeur paramétrise
$q_\phi(\mathbf{z}|\mathbf{x}) =
\mathcal{N}(\mathbf{z};\,\boldsymbol{\mu}_\phi(\mathbf{x}),\,
\mathrm{diag}(\boldsymbol{\sigma}^2_\phi(\mathbf{x})))$ via deux têtes
linéaires distinctes. L'astuce du reparamétrage rend l'échantillonnage
différentiable :

$$\mathbf{z} = \boldsymbol{\mu}_\phi(\mathbf{x}) +
\boldsymbol{\sigma}_\phi(\mathbf{x}) \odot \boldsymbol{\epsilon},
\qquad \boldsymbol{\epsilon} \sim \mathcal{N}(\mathbf{0}, \mathbf{I})$$

\textbf{Terme KL en forme close.} Pour des gaussiennes diagonales :

$$D_{\mathrm{KL}}\!\left(\mathcal{N}(\boldsymbol{\mu},\boldsymbol{\sigma}^2)
\;\|\; \mathcal{N}(\mathbf{0},\mathbf{I})\right) =
-\frac{1}{2}\sum_{j=1}^{d}\left(1 + \log\sigma_j^2 - \mu_j^2 -
\sigma_j^2\right)$$

\textbf{Score d'anomalie dérivé (An \& Cho, 2015) :}

$$\mathrm{AnoScore}(\mathbf{x}) = -\frac{1}{L}\sum_{\ell=1}^{L}
\log p_\theta\!\left(\mathbf{x}\mid\mathbf{z}^{(\ell)}\right),
\quad \mathbf{z}^{(\ell)} \sim q_\phi(\mathbf{z}|\mathbf{x})$$

\textbf{Limites structurelles.} (1)~Indépendance temporelle : chaque observation
est traitée indépendamment. (2)~Prior gaussienne standard limitant l'expressivité
face aux données financières hétéroscédastiques. (3)~Factorisation diagonale de
$q_\phi$ ignorant les corrélations entre dimensions latentes.

\textbf{Positionnement benchmark.} Travail fondationnel sans évaluation sur les
datasets SMD/MSL/SMAP/SWaT/PSM : le modèle ne traite pas de séries temporelles.
Il constitue la brique de base de tous les modèles évalués dans le tableau de
la section~7.

\textbf{Critères architecturaux.} Encodeur non causal (statique) ; score
probabiliste (reconstruction probability) ; explicabilité nulle ; données
tabulaires génériques ; complexité attention sans objet.

\textbf{Pertinence pour BiVAT.} Fondamentale. L'ELBO, le reparamétrage et le
score par reconstruction probability constituent les briques de base de BiVAT.

% ──────────────────────────────────────────────────────────────
\subsection{[F2] An \& Cho (2015)  VAE-Based Anomaly Detection Using
Reconstruction Probability}

\textbf{Problème formel.} Comment formaliser un critère de décision probabiliste
exploitant la nature stochastique de l'espace latent pour quantifier l'anomalie
au-delà de la simple erreur quadratique de reconstruction ?

\textbf{Contribution principale.} An \& Cho formalisent la \textit{reconstruction
probability} comme score d'anomalie :

$$p_{\mathrm{rec}}(\mathbf{x}) = \frac{1}{L}\sum_{\ell=1}^{L}
\log p_\theta\!\left(\mathbf{x}\mid\mathbf{z}^{(\ell)}\right),
\quad \mathbf{z}^{(\ell)} \sim q_\phi(\mathbf{z}|\mathbf{x})$$

Pour $p_\theta(\mathbf{x}|\mathbf{z}) = \mathcal{N}(\hat{\mathbf{x}},
\sigma^2\mathbf{I})$ :

$$p_{\mathrm{rec}}(\mathbf{x}) \approx
-\frac{1}{2\sigma^2 L}\sum_{\ell=1}^{L}
\left\|\mathbf{x} - \hat{\mathbf{x}}^{(\ell)}\right\|^2 + \mathrm{const}$$

Les $L$ tirages Monte Carlo permettent d'estimer l'incertitude de reconstruction :
pour un point atypique, les reconstructions $\hat{\mathbf{x}}^{(\ell)}$ divergent
fortement entre les tirages, amplifiant le signal d'anomalie. La règle de décision
calibre un seuil $\tau$ au $(1-\alpha)$-ième percentile empirique.

\textbf{Limites.} Application uniquement sur données statiques ; absence de
modélisation des dépendances temporelles ; coût Monte Carlo à l'inférence.

\textbf{Positionnement benchmark.} Travail fondationnel sur données statiques :
aucun résultat sur SMD/MSL/SMAP/SWaT/PSM. La reconstruction probability est
reprise comme composante du score dans OmniAnomaly, T-VAE et BiVAT.

\textbf{Critères architecturaux.} Encodeur non causal (statique) ; score
probabiliste ($L$ tirages MC) ; explicabilité partielle (variance entre tirages) ;
multivarié natif ; complexité attention sans objet ; domaine médical et générique.

\textbf{Pertinence pour BiVAT.} Élevée. La reconstruction probability constitue
le composant VAE du score hybride BiVAT.

% ──────────────────────────────────────────────────────────────
\subsection{[F3] Su et al. (2019, KDD)  OmniAnomaly}

\textbf{Problème formel.} Détecter des anomalies dans des séries temporelles
multivariées $\mathbf{X} \in \mathbb{R}^{T \times D}$ en tenant compte
simultanément des dépendances temporelles séquentielles et des corrélations
inter-dimensions.

\textbf{Architecture : GRU-VAE stochastique.}

\textit{Encodeur :}
$$\mathbf{h}^e_t = \mathrm{GRU}_\phi^e(\mathbf{x}_t,\,\mathbf{h}^e_{t-1}),
\qquad \mathbf{z}_t = \boldsymbol{\mu}^e_t +
\boldsymbol{\sigma}^e_t \odot \boldsymbol{\epsilon}_t$$

\textit{Prior temporelle conditionnelle (innovation centrale) :}
$$p(\mathbf{z}_t | \mathbf{z}_{t-1}) = \mathcal{N}\!\left(\mathbf{z}_t;\;
\boldsymbol{\mu}^p_t(\mathbf{h}^d_{t-1}),\;
\boldsymbol{\sigma}^p_t(\mathbf{h}^d_{t-1})\right)$$

\textbf{Fonction objectif (ELBO séquentiel) :}
$$\mathcal{L}_{\mathrm{OmniAnomaly}} = \sum_{t=1}^T\!\left[
\mathbb{E}_{q_\phi}\!\left[\log p_\theta(\mathbf{x}_t|\mathbf{z}_t)\right]
- D_{\mathrm{KL}}\!\left(q_\phi(\mathbf{z}_t|\mathbf{x}_{\leq t})\;\|\;
p(\mathbf{z}_t|\mathbf{z}_{t-1})\right)\right]$$

\textbf{Limites.} Fenêtre de contexte limitée (oubli GRU pour $T > 50$) ;
traitement séquentiel interdisant la parallélisation GPU ; absence d'attention
globale.

\textbf{Positionnement benchmark.}

\smallskip
\begin{center}
\small
\begin{tabular}{lccccc}
\toprule
& SMD & MSL & SMAP & SWaT & PSM \\
\midrule
P (\%) & 83.68 & 89.02 & 92.49 & 81.42 & 88.39 \\
R (\%) & 86.82 & 86.37 & 81.99 & 84.30 & 74.46 \\
F1 (\%) & 85.22 & 87.67 & 86.92 & 82.83 & 80.83 \\
\bottomrule
\end{tabular}
\end{center}
\smallskip
\noindent\textit{\small Source : Xu et al.\ 2022, Table~1 (point adjustment).
F1 moyen sur 5 datasets : 84.69\%.}

\textbf{Critères architecturaux.} Encodeur bidirectionnel (LSTM stochastique +
normalizing flow) ; score probabiliste ($K$ tirages MC) ; explicabilité partielle
(par dimension) ; multivarié natif (38 dimensions SMD) ;
complexité $\mathcal{O}(T)$ LSTM ; domaine industriel/capteurs.

\textbf{Pertinence pour BiVAT.} Élevée. La prior temporelle conditionnelle et
le score probabiliste séquentiel sont réutilisables. Le GRU sera remplacé par
le module hybride Transformer + BiLSTM de BiVAT.

% ════════════════════════════════════════════════════════════════
\section{Axe 2 : Transformers pour la Détection d'Anomalies}

\subsection{[T1] Xu et al. (2022, ICLR Spotlight)  Anomaly Transformer}

\textbf{Problème formel.} La détection non supervisée de points anomaux dans
les séries temporelles exige du modèle qu'il dérive un critère distinctif
exploitant la structure temporelle globale, au-delà des associations point-à-point.

\textbf{Observation fondatrice.} En raison de la rareté des anomalies, il est
extrêmement difficile de construire des associations non triviales depuis les
points anomaux vers l'ensemble de la série. Les associations des anomalies se
concentrent donc principalement sur leurs points temporels adjacents. Ce biais
de concentration adjacente implique un critère discriminant inhérent entre
points normaux et anormaux.

\textbf{Association Discrepancy  Formalisation.}

\textit{Prior Association :}
$$\mathcal{P}^{(\ell)}_t(i) = \frac{1}{\sqrt{2\pi}\sigma^{(\ell)}}
\exp\!\left(-\frac{(i-t)^2}{2(\sigma^{(\ell)})^2}\right),
\quad i \in \{1,\ldots,T\}$$

\textit{Series Association :}
$$\mathcal{S}^{(\ell)} = \mathrm{Softmax}\!\left(
\frac{\mathbf{Q}^{(\ell)}{\mathbf{K}^{(\ell)}}^\top}{\sqrt{d_k}}\right)$$

\textit{Association Discrepancy (divergence KL symétrique) :}
$$\mathrm{AssDis}^{(\ell)}(\mathbf{X}) = \frac{1}{T}\sum_{t=1}^{T}\!\left[
\mathrm{KL}\!\left(\mathcal{P}^{(\ell)}_t \;\|\; \mathcal{S}^{(\ell)}_t\right)
+ \mathrm{KL}\!\left(\mathcal{S}^{(\ell)}_t \;\|\;
\mathcal{P}^{(\ell)}_t\right)\right]$$

\textbf{Stratégie d'entraînement minimax.}

\textit{Phase Maximize :}
$$\min_{\theta,\sigma}\left\|\mathbf{X} - \hat{\mathbf{X}}\right\|_F^2
- \lambda\cdot\mathrm{AssDis}(\mathbf{X})$$

\textit{Phase Minimize :}
$$\min_{\theta,\sigma}\left\|\mathbf{X} - \hat{\mathbf{X}}\right\|_F^2
+ \lambda\cdot\mathrm{AssDis}(\mathbf{X})$$

\textbf{Score d'anomalie final :}
$$\mathrm{AnoScore}(t) = \mathrm{Softmax}\!\left(
-\mathrm{AssDis}^{(1:L)}_t\right)\cdot
\left\|\mathbf{x}_t - \hat{\mathbf{x}}_t\right\|_2$$

\textbf{Architecture.} Stack de $L = 3$ blocs Transformer modifiés avec
embedding positionnel sinusoïdal. Complexité $\mathcal{O}(T^2 d)$.

\textbf{Limites.} Score déterministe sans intervalle de confiance ; encodeur
causal uniquement ; score scalaire unique sans distinction des types d'anomalies.

\textbf{Positionnement benchmark.}

\smallskip
\begin{center}
\small
\begin{tabular}{lccccc}
\toprule
& SMD & MSL & SMAP & SWaT & PSM \\
\midrule
P (\%) & 89.40 & 92.09 & 94.13 & 91.55 & 96.91 \\
R (\%) & 95.45 & 95.15 & 99.40 & 96.73 & 98.90 \\
F1 (\%) & 92.33 & 93.59 & 96.69 & 94.07 & 97.89 \\
\bottomrule
\end{tabular}
\end{center}
\smallskip
\noindent\textit{\small Source : Xu et al.\ 2022, Table~1 (point adjustment).
F1 moyen sur 5 datasets : 94.91\%. Modèle de référence dominant sur le panel
normalisé.}

\textbf{Critères architecturaux.} Encodeur causal (minimax + prior gaussienne
locale) ; score déterministe ; explicabilité partielle (attention weights) ;
multivarié natif (55 dimensions MSL) ; complexité $\mathcal{O}(T^2)$ ;
domaine industriel, spatial, eau.

\textbf{Pertinence pour BiVAT.} Très élevée. L'Association Discrepancy et la
stratégie minimax sont intégrées dans BiVAT. Le score hybride reconstruction
$\times$ AssDis apporte $+18{,}82$ points F1 moyen par rapport à la
reconstruction seule (ablation study complète en section~8.4).

% ──────────────────────────────────────────────────────────────
\subsection{[T2] Tuli, Casale \& Jennings (2022, VLDB)  TranAD}

\textbf{Problème formel.} Réaliser simultanément la détection et le diagnostic
d'anomalies dans des données temporelles multivariées dans un cadre Transformer
unique, à inférence rapide, robuste face à la multi-modalité des données.

\textbf{Architecture : deux décodeurs antagonistes.}

\textit{Encodeur partagé :}
$$\mathbf{H} = \mathrm{TransformerEncoder}(\mathbf{X}_{t-W:t})
\in \mathbb{R}^{W \times d_{\mathrm{model}}}$$

\textit{Décodeur 1 et focus score :}
$$\hat{\mathbf{X}}^{(1)} = \mathcal{D}_1(\mathbf{H}), \qquad
\mathbf{f}^{(1)}_t = \left\|\mathbf{x}_t -
\hat{\mathbf{x}}^{(1)}_t\right\|_1 \in \mathbb{R}^D$$

\textit{Décodeur 2 conditionné :}
$$\hat{\mathbf{X}}^{(2)} = \mathcal{D}_2(\mathbf{H},\,\mathbf{f}^{(1)})$$

\textbf{Fonction objectif adversariale.} Avec $n$ l'indice d'époque :
$$\mathcal{L} = \frac{1}{n}\,\mathcal{L}^{(1)} +
\left(1 - \frac{1}{n}\right)\mathcal{L}^{(2)},
\qquad \mathcal{L}^{(i)} = \frac{1}{W}\sum_t
\|\mathbf{x}_t - \hat{\mathbf{x}}^{(i)}_t\|_2^2$$

\textbf{Avantage computationnel documenté.} TranAD réduit les temps
d'entraînement de 75 à 99\% par rapport aux baselines BiLSTM/LSTM grâce à la
parallélisation du Transformer (Tuli et al.\ 2022, Table~5).

\textbf{Limites.} Absence de composante variationnelle ; score déterministe
sans intervalle de confiance ; pas d'explicabilité formelle.

\textbf{Positionnement benchmark.}

\smallskip
\begin{center}
\small
\begin{tabular}{lccccc}
\toprule
& SMD & MSL & SMAP & SWaT & PSM \\
\midrule
P (\%) & 92.62 & 90.38 & 80.43 & 97.60 & \nd \\
R (\%) & 99.74 & 99.99 & 99.99 & 69.97 & \nd \\
F1 (\%) & 96.05 & 94.94 & 89.15 & 81.51 & \nd \\
\bottomrule
\end{tabular}
\end{center}
\smallskip
\noindent\textit{\small Source : Tuli et al.\ 2022, Table~2 (point adjustment).
PSM non évalué. F1 moyen sur 4 datasets : 90.41\%. Meilleur F1 sur SMD (96.05\%).}

\textbf{Critères architecturaux.} Encodeur non causal (attention bidirectionnelle
+ focus score) ; score déterministe (deux phases) ; explicabilité partielle
(focus scores par dimension) ; multivarié natif (123 dimensions WADI) ;
complexité $\mathcal{O}(T^2)$ ; domaine industriel, IoT.

\textbf{Pertinence pour BiVAT.} Élevée. Le mécanisme de focus score est
transposable comme signal de conditionnement entre les composantes VAE et
Transformer de l'encodeur hybride BiVAT.

% ════════════════════════════════════════════════════════════════
\section{Axe 3 : Hybrides VAE-Transformer}

\subsection{[H1] Song, Seo \& Kim (2023, IEEE Access) 
Anomaly VAE-Transformer}

\textbf{Problème formel.} Les séries temporelles financières (protocoles DeFi)
présentent des anomalies à deux échelles distinctes : des pics ponctuels
détectables localement par le VAE, et des dérives progressives détectables sur
de longs horizons via le Transformer.

\textbf{Architecture à double branche.}

$$s^{\mathrm{VAE}}_t = -\log p_\theta(\mathbf{w}_t|\mathbf{z}_t), \qquad
s^{\mathrm{Tr}}_t = \|\mathbf{x}_t - \hat{\mathbf{x}}^{\mathrm{Tr}}_t\|_2$$

$$s_t^{\mathrm{final}} = \alpha\,s^{\mathrm{VAE}}_t +
(1-\alpha)\,s^{\mathrm{Tr}}_t, \quad \alpha \in [0,1]\text{ appris}$$

\textbf{Résultats sur données DeFi.} F1 = 91.3\% contre 87.2\% pour le VAE
seul et 88.9\% pour le Transformer seul (Olympus DAO, Uniswap, Aave). Ce modèle
est le seul de la liste validé sur des données financières réelles.

\textbf{Limite fondamentale.} L'intégration est superficielle : VAE et
Transformer opèrent dans des espaces distincts, leurs gradients ne s'informent
pas mutuellement. La fusion est une pondération de scalaires, non une
intégration dans l'espace de représentation.

\textbf{Positionnement benchmark.} Non évalué sur SMD, MSL, SMAP, SWaT, PSM.
Données DeFi exclusivement ; comparaison normalisée impossible sans
réimplémentation. Maintenu dans la matrice qualitative uniquement.

\textbf{Critères architecturaux.} Encodeur non causal (VAE local + Transformer
long-terme) ; score hybride (scalaire VAE + Transformer) ; pas d'explicabilité ;
multivarié partiel (DeFi uni/bi-varié) ; complexité $\mathcal{O}(T^2)$ ;
\textbf{seul modèle validé sur données financières réelles (DeFi blockchain)}.

\textbf{Pertinence pour BiVAT.} Modérée. Ce papier valide la complémentarité
VAE + Transformer sur données financières, mais l'intégration superficielle est
la limite précise que BiVAT surmonte par une intégration profonde dans l'espace
de représentation.

% ──────────────────────────────────────────────────────────────
\subsection{[H2] Li et al. (2025, Expert Systems)  T-VAE}

\textbf{Problème formel.} La détection d'anomalies dans des séries temporelles
multivariées dont les dépendances séquentielles et inter-dimensions sont
complexes, non linéaires, avec forte volatilité.

\textbf{Architecture : intégration profonde Transformer-dans-VAE.}

\textit{Étape 1  Embedding :}
$\mathbf{E}_t = \mathrm{EmbedLayer}(\mathbf{x}_t) + \mathrm{PE}(t)$

\textit{Étape 2  Transformer ($L$ couches) :}
$\mathbf{H}^{(L)} = \mathrm{TransformerEncoder}^{(1:L)}(\mathbf{E})$

\textit{Étape 3  Projection variationnelle :}
$$\boldsymbol{\mu}_t = \mathbf{W}_\mu\,\mathbf{H}^{(L)}_t + \mathbf{b}_\mu,
\qquad \log\boldsymbol{\sigma}^2_t = \mathbf{W}_\sigma\,\mathbf{H}^{(L)}_t
+ \mathbf{b}_\sigma$$

\textbf{Fonction objectif avec KL Annealing :}
$$\mathcal{L}_{\mathrm{T-VAE}} = \sum_{t=1}^T\left[
\mathbb{E}_{q_\phi}\!\left[\log p_\theta(\mathbf{x}_t|\mathbf{z}_t)\right]
- \beta_t\,D_{\mathrm{KL}}\!\left(q_\phi(\mathbf{z}_t|\mathbf{x})
\;\|\; \mathcal{N}(\mathbf{0},\mathbf{I})\right)\right]$$

\textbf{Score d'anomalie :}
$$s_t = -\frac{1}{L}\sum_{\ell=1}^L\log p_\theta\!\left(
\mathbf{x}_t\mid\mathbf{z}^{(\ell)}_t\right) +
\gamma\,D_{\mathrm{KL}}(q_\phi(\mathbf{z}_t|\mathbf{x})\;\|\;
\mathcal{N}(\mathbf{0},\mathbf{I}))$$

\textbf{Limites identifiées.} Encodeur causal (attention masquée, pas
d'information future dans l'espace latent) ; score scalaire unique ; absence
d'explicabilité formelle.

\textbf{Positionnement benchmark.} Derrière paywall Wiley. Tables P/R/F1 sur
les 5 datasets de référence non accessibles en libre accès. Aucune
réimplémentation afin de préserver la rigueur de la comparaison. Maintenu
dans la matrice qualitative en tant que référence architecturale directe.

\textbf{Critères architecturaux.} Encodeur bidirectionnel (Transformer +
projection VAE) ; score probabiliste ; pas d'explicabilité ; multivarié natif ;
complexité $\mathcal{O}(T^2)$ ; données industrielles.

\textbf{Pertinence pour BiVAT.} Très élevée. T-VAE est l'architecture la plus
proche de BiVAT. Les contributions de BiVAT par rapport au T-VAE sont :
(i)~encodeur hybride Transformer + BiLSTM pour les petits volumes ; (ii)~Conformal
Prediction pour les intervalles de confiance ; (iii)~Association Discrepancy
dans la fonction objectif.

% ──────────────────────────────────────────────────────────────
\subsection{[H3] Wang et al. (2022, Measurement) 
Variational Transformer-based Anomaly Detection}

\textbf{Problème formel.} Formaliser l'intégration de la stochasticité
variationnelle au sein du flux d'attention du Transformer.

\textbf{Architecture : couche d'attention variationnelle.} Après le calcul de
$\mathrm{Softmax}(\mathbf{QK}^\top/\sqrt{d_k})$, un MLP produit
$\boldsymbol{\mu}_{\mathrm{attn}}$ et $\boldsymbol{\sigma}_{\mathrm{attn}}$ :
$$\mathbf{z}_{\mathrm{attn}} \sim \mathcal{N}(\boldsymbol{\mu}_{\mathrm{attn}},
\,\mathrm{diag}(\boldsymbol{\sigma}^2_{\mathrm{attn}})), \qquad
\mathbf{H} = \mathbf{z}_{\mathrm{attn}}\,\mathbf{V}$$

\textbf{Score d'anomalie :}
$$s_t = -\log p_\theta(\mathbf{x}_t|\mathbf{z}_{\mathrm{attn},\leq t})
+ \lambda\,D_{\mathrm{KL}}(\mathcal{N}(\boldsymbol{\mu}_{\mathrm{attn},t},
\boldsymbol{\sigma}^2_{\mathrm{attn},t})\|\mathcal{N}(\mathbf{0},\mathbf{I}))$$

\textbf{Positionnement benchmark.} Derrière paywall Elsevier. Résultats non
accessibles. Maintenu dans la matrice qualitative uniquement.

\textbf{Critères architecturaux.} Encodeur non causal (Transformer + VAE sur
espace d'attention) ; score probabiliste ; pas d'explicabilité formelle ;
multivarié natif ; complexité $\mathcal{O}(T^2)$ ; données industrielles.

\textbf{Pertinence pour BiVAT.} Élevée. L'injection de stochasticité dans
le mécanisme d'attention ouvre la voie à des distributions d'attention
incertaines interprétables comme des distributions sur les associations
temporelles.

% ──────────────────────────────────────────────────────────────
\subsection{[H4] Zhang et al. (2021, CCDC) 
Transformer-based VAE Non Supervisé}

\textbf{Positionnement historique.} Premier travail proposant explicitement le
remplacement des couches GRU d'OmniAnomaly par des blocs d'attention multi-têtes,
constituant le premier Transformer-VAE pour la détection d'anomalies non
supervisée dans les séries multivariées.

\textbf{Architecture.}
$$\mathbf{H} = \mathrm{MultiHeadAttn}(\mathbf{X}, \mathbf{X}, \mathbf{X})
+ \mathbf{X}, \qquad
\boldsymbol{\mu}_t,\,\boldsymbol{\sigma}_t = \mathrm{MLP}(\mathbf{H}_t)$$

Le décodeur reste un LSTM. Gains rapportés : $+2$ à $+5$ points F1 sur
OmniAnomaly sur SMD et MSL.

\textbf{Positionnement benchmark.} Résultats disponibles sur SMD et MSL
uniquement. Non inclus dans le tableau normalisé à 5 datasets.

\textbf{Critères architecturaux.} Encodeur bidirectionnel (Transformer non
masqué) ; score probabiliste (VAE) ; pas d'explicabilité ; multivarié natif ;
complexité $\mathcal{O}(T^2)$ ; données industrielles.

\textbf{Pertinence pour BiVAT.} Élevée. Pont historique essentiel pour la
narration OmniAnomaly $\rightarrow$ Transformer-VAE $\rightarrow$ BiVAT.

% ════════════════════════════════════════════════════════════════
\section{Axe 4 : Architectures Avancées}

\subsection{[A1] He et al. (2024, Information Sciences)  VAEAT}

\textbf{Problème formel.} Les données d'entraînement d'un VAE peuvent être
contaminées par des anomalies non étiquetées, dégradant la qualité de
reconstruction. Comment rendre le VAE robuste via un entraînement adversarial ?

\textbf{Architecture : deux VAE hybrides LSTM + Attention.}

$$\mathbf{h}^{\mathrm{LSTM}}_t =
\mathrm{LSTM}_\phi(\mathbf{x}_t,\,\mathbf{h}^{\mathrm{LSTM}}_{t-1}), \qquad
\mathbf{h}^{\mathrm{attn}}_t =
\mathrm{MHA}(\mathbf{h}^{\mathrm{LSTM}},
\mathbf{h}^{\mathrm{LSTM}},\mathbf{h}^{\mathrm{LSTM}})_t$$
$$\tilde{\mathbf{h}}_t = \mathrm{LayerNorm}(\mathbf{h}^{\mathrm{LSTM}}_t
+ \mathbf{h}^{\mathrm{attn}}_t)$$

\textbf{Stratégie en deux phases.} Phase~1 : deux VAE reconstituent les
données. Phase~2 : un discriminateur $\mathcal{D}$ distingue les résidus
normaux des résidus anormaux.

$$\mathcal{L}_{\mathrm{VAEAT}} = \mathcal{L}^{(1)}_{\mathrm{ELBO}}
+ \mathcal{L}^{(2)}_{\mathrm{ELBO}} + \lambda\,\mathcal{L}_{\mathrm{adv}},
\qquad
s_t = \alpha\,s^{(1)}_t + (1-\alpha)\,s^{(2)}_t +
\beta\,(1-\mathcal{D}(\mathbf{r}_t))$$

\textbf{Positionnement benchmark.}

\smallskip
\begin{center}
\small
\begin{tabular}{lccccc}
\toprule
& SMD & MSL & SMAP & SWaT & PSM \\
\midrule
P (\%) & 95.18 & 95.22 & 89.43 & 93.63 & 95.78 \\
R (\%) & 96.54 & 99.99 & 99.51 & 93.50 & 96.74 \\
F1 (\%) & 95.71 & 97.35 & 92.96 & 93.57 & 96.26 \\
\bottomrule
\end{tabular}
\end{center}
\smallskip
\noindent\textit{\small Source : He et al.\ 2024, Table~3 (point adjustment).
F1 moyen : 95.17\%. Meilleurs F1 VAE sur SMD et MSL du panel normalisé.}

\textbf{Critères architecturaux.} Encodeur non causal (LSTM + attention
bidirectionnelle) ; score probabiliste (VAE + discriminateur) ; pas
d'explicabilité ; multivarié natif (55 dimensions MSL) ; complexité $\mathcal{O}(T)$
LSTM ; données industrielles.

\textbf{Pertinence pour BiVAT.} Très élevée. VAEAT valide empiriquement que
l'hybride LSTM + attention bidirectionnelle atteint les meilleures performances
VAE sur SMD (95.71\%) et MSL (97.35\%). La stratégie adversariale deux phases
est adaptable pour les chocs COVID-19 non étiquetés du dataset BEAC.

% ──────────────────────────────────────────────────────────────
\subsection{[A2] Xie et al. (2025, CMC)  VLT-Anomaly}

\textbf{Problème formel.} Adresser simultanément l'absence de dépendances
bidirectionnelles, les dépendances long-terme insuffisantes et la détection à
grain unique des méthodes existantes.

\textbf{Solutions architecturales.} $\beta$-VAE avec hyperparamètre KL
contrôlé ; module hybride Transformer + BiLSTM (le Transformer pour les
dépendances globales, le BiLSTM pour la modélisation bidirectionnelle locale) ;
score par moyennage prédiction-reconstruction. L'ablation (Table~3, Xie et al.\
2025) confirme que BiLSTM seul, Transformer seul et VAE seul sont tous
inférieurs à l'hybride complet.

\textbf{Positionnement benchmark.} Évalué exclusivement sur NAB (séries
univariées). Aucun résultat sur SMD/MSL/SMAP/SWaT/PSM. Maintenu dans la
matrice qualitative ; cité pour son ablation study sur l'hybride
Transformer + BiLSTM.

\textbf{Critères architecturaux.} Encodeur bidirectionnel ($\beta$-VAE +
BiLSTM + Transformer) ; score déterministe ; pas d'explicabilité ; univarié
(NAB) ; complexité $\mathcal{O}(T^2) + \mathcal{O}(T)$ ; données partiellement
financières (NAB).

\textbf{Pertinence pour BiVAT.} Très élevée. VLT-Anomaly fournit la validation
de principe la plus directe de l'encodeur hybride Transformer + BiLSTM retenu
dans BiVAT. L'ablation Table~3 est citée en section~8.1.

% ──────────────────────────────────────────────────────────────
\subsection{[A3] Zhang et al. (2025, Engineering Applications of AI) 
SHAPAttenVAE}

\textbf{Problème formel.} Rendre explicable chaque décision d'anomalie d'un VAE
par attribution quantitative dimensionnelle, dans un cadre réglementaire exigeant
l'auditabilité.

\textbf{Architecture et mécanisme d'explicabilité.}

\textit{Poids d'attention inter-couches :}
$$\boldsymbol{\alpha}^{(\ell)}_t =
\mathrm{Softmax}(\mathbf{W}^{(\ell)}\,\mathbf{h}^{(\ell-1)}_t)
\in \mathbb{R}^D$$

\textit{Valeurs SHAP locales :}
$$\phi_d(\mathbf{x}_t) = \sum_{S \subseteq \mathcal{D}\setminus\{d\}}
\frac{|S|!\,(D-|S|-1)!}{D!}
\left[s(\mathbf{x}_{S\cup\{d\}}) - s(\mathbf{x}_S)\right]$$

\textit{Perte de cohérence :}
$$\mathcal{L}_{\mathrm{expl}} = \sum_t
\left\|\boldsymbol{\alpha}_t - \hat{\boldsymbol{\phi}}_t\right\|_2^2$$

\textbf{Positionnement benchmark.} Derrière paywall ScienceDirect
(DOI: 10.1016/j.engappai.2025.111585). Résultats non accessibles. Maintenu
dans la matrice qualitative ; mécanisme SHAP inter-couches directement repris
dans BiVAT.

\textbf{Critères architecturaux.} Encodeur non causal (MHA + VAE) ; score
probabiliste ; \textbf{explicabilité SHAP inter-couches formelle} ; multivarié
natif ; complexité $\mathcal{O}(T^2)$ ; données industrielles.

\textbf{Pertinence pour BiVAT.} Très élevée. SHAPAttenVAE est la seule
architecture de la littérature à avoir intégré SHAP formellement dans un VAE
avec couches d'attention. L'exigence COBAC fait de ce mécanisme une nécessité
fonctionnelle.

% ──────────────────────────────────────────────────────────────
\subsection{[A4] Sellam et al. (2025, Engineering Applications of AI)  MAAT}

\textbf{Problème formel.} Conserver la qualité de modélisation des dépendances
longue portée tout en réduisant la complexité $\mathcal{O}(T^2)$ de l'attention
standard à $\mathcal{O}(T)$.

\textbf{Mamba : State Space Model sélectif.}
$$\mathbf{h}_t = \bar{\mathbf{A}}\,\mathbf{h}_{t-1} +
\bar{\mathbf{B}}\,\mathbf{u}_t, \qquad \mathbf{y}_t = \mathbf{C}\,\mathbf{h}_t$$

\textbf{Architecture MAAT.} Sparse Attention $\mathcal{O}(T \log T)$ +
Mamba-SSM $\mathcal{O}(T)$ + Gated Attention adaptatif.

\textbf{Positionnement benchmark.}

\smallskip
\begin{center}
\small
\begin{tabular}{lccccc}
\toprule
& SMD & MSL & SMAP & SWaT & PSM \\
\midrule
P (\%) & 90.03 & 93.23 & 95.37 & 91.27 & \nd \\
R (\%) & 99.29 & 97.55 & 99.32 & 97.32 & \nd \\
F1 (\%) & 94.51 & 95.05 & 96.99 & 94.16 & 98.76 \\
\bottomrule
\end{tabular}
\end{center}
\smallskip
\noindent\textit{\small Source : arXiv:2502.07858, Table~1 (deltas par rapport
à l'Anomaly Transformer). F1 moyen : 95.89\%. Modèle le plus performant du
panel normalisé.}

\textbf{Critères architecturaux.} Encodeur causal + Mamba SSM ; score
déterministe (AssDis) ; pas d'explicabilité ; multivarié natif (38 dimensions) ;
complexité $\mathcal{O}(T)$ Mamba / $\mathcal{O}(T \log T)$ sparse attention ;
données industrielles.

\textbf{Pertinence pour BiVAT.} Modérée. Pour $T = 195$, la complexité
quadratique est négligeable. Mamba constitue une variante d'encodeur intégrable
si BiVAT est déployé sur fréquences journalières ou hebdomadaires.

% ════════════════════════════════════════════════════════════════
\section{Synthèse Transversale et Cartographie des Lacunes}

\subsection{Tableau de benchmarks normalisés}

\textit{Note méthodologique.} Toutes les métriques sont en pourcentage. Le
protocole majoritaire est le \textit{point adjustment} (PA) : si au moins un
point d'un segment anormal est détecté, tout le segment est crédité. Ce
protocole tend à gonfler les scores par rapport à une évaluation strictement
point-à-point. Les cinq modèles marqués \nd{} (T-VAE, VLT-Anomaly,
SHAPAttenVAE, Anomaly VAE-Transformer, Variational Transformer) ne disposent
pas de résultats publics sur ce panel : soit les papiers sont derrière paywall
(T-VAE, SHAPAttenVAE, Variational Transformer), soit les modèles ont été
évalués sur des datasets incompatibles (VLT-Anomaly sur NAB, Anomaly
VAE-Transformer sur données DeFi). Aucune réimplémentation n'a été réalisée
afin de préserver la rigueur de la comparaison.

\begin{landscape}
\begin{longtable}{llcccl}
\caption{Benchmarks normalisés  Precision (P), Recall (R), F1 (\%).}
\label{tab:benchmarks}\\
\toprule
\textbf{Modèle} & \textbf{Dataset} & \textbf{P} & \textbf{R} &
\textbf{F1} & \textbf{Source} \\
\midrule
\endfirsthead
\multicolumn{6}{c}{\tablename~\thetable{}  Suite}\\
\toprule
\textbf{Modèle} & \textbf{Dataset} & \textbf{P} & \textbf{R} &
\textbf{F1} & \textbf{Source} \\
\midrule
\endhead
\midrule\multicolumn{6}{r}{Suite page suivante\ldots}\endfoot
\bottomrule\endlastfoot

\multirow{5}{*}{\textbf{OmniAnomaly}}
 & SMD  & 83.68 & 86.82 & 85.22 & Xu et al.\ 2022, Table~1 \\
 & MSL  & 89.02 & 86.37 & 87.67 & idem \\
 & SMAP & 92.49 & 81.99 & 86.92 & idem \\
 & SWaT & 81.42 & 84.30 & 82.83 & idem \\
 & PSM  & 88.39 & 74.46 & 80.83 & idem \\
\midrule
\multirow{5}{*}{\textbf{Anomaly Transformer}}
 & SMD  & 89.40 & 95.45 & 92.33 & Xu et al.\ 2022, Table~1 \\
 & MSL  & 92.09 & 95.15 & 93.59 & idem \\
 & SMAP & 94.13 & 99.40 & 96.69 & idem \\
 & SWaT & 91.55 & 96.73 & 94.07 & idem \\
 & PSM  & 96.91 & 98.90 & 97.89 & idem \\
\midrule
\multirow{5}{*}{\textbf{TranAD}}
 & SMD  & 92.62 & 99.74 & 96.05 & Tuli et al.\ 2022, Table~2 \\
 & MSL  & 90.38 & 99.99 & 94.94 & idem \\
 & SMAP & 80.43 & 99.99 & 89.15 & idem \\
 & SWaT & 97.60 & 69.97 & 81.51 & idem \\
 & PSM  & \nd   & \nd   & \nd   & Non évalué dans le papier \\
\midrule
\multirow{5}{*}{\textbf{T-VAE}}
 & SMD  & \nd & \nd & \nd &
   \multirow{5}{*}{Li et al.\ 2025  paywall Wiley} \\
 & MSL  & \nd & \nd & \nd & \\
 & SMAP & \nd & \nd & \nd & \\
 & SWaT & \nd & \nd & \nd & \\
 & PSM  & \nd & \nd & \nd & \\
\midrule
\multirow{5}{*}{\textbf{VLT-Anomaly}}
 & SMD  & \nd & \nd & \nd &
   \multirow{5}{*}{Xie et al.\ 2025  NAB uniquement} \\
 & MSL  & \nd & \nd & \nd & \\
 & SMAP & \nd & \nd & \nd & \\
 & SWaT & \nd & \nd & \nd & \\
 & PSM  & \nd & \nd & \nd & \\
\midrule
\multirow{5}{*}{\textbf{SHAPAttenVAE}}
 & SMD  & \nd & \nd & \nd &
   \multirow{5}{*}{Zhang et al.\ 2025  paywall ScienceDirect} \\
 & MSL  & \nd & \nd & \nd & \\
 & SMAP & \nd & \nd & \nd & \\
 & SWaT & \nd & \nd & \nd & \\
 & PSM  & \nd & \nd & \nd & \\
\midrule
\multirow{5}{*}{\textbf{MAAT}}
 & SMD  & 90.03 & 99.29 & 94.51 & arXiv:2502.07858, Table~1 \\
 & MSL  & 93.23 & 97.55 & 95.05 & idem ($+1.46$ pts vs AT) \\
 & SMAP & 95.37 & 99.32 & 96.99 & idem ($+0.30$ pts vs AT) \\
 & SWaT & 91.27 & 97.32 & 94.16 & idem ($+0.09$ pts vs AT) \\
 & PSM  & \nd   & \nd   & 98.76 & idem ($+0.87$ pts vs AT) \\
\midrule
\multirow{5}{*}{\textbf{VAEAT}}
 & SMD  & 95.18 & 96.54 & 95.71 & He et al.\ 2024, Table~3 \\
 & MSL  & 95.22 & 99.99 & 97.35 & idem \\
 & SMAP & 89.43 & 99.51 & 92.96 & idem \\
 & SWaT & 93.63 & 93.50 & 93.57 & idem \\
 & PSM  & 95.78 & 96.74 & 96.26 & idem \\
\midrule
\multirow{5}{*}{\textbf{Anomaly VAE-Transf.}}
 & SMD  & \nd & \nd & \nd &
   \multirow{5}{*}{Song et al.\ 2023  données DeFi uniquement} \\
 & MSL  & \nd & \nd & \nd & \\
 & SMAP & \nd & \nd & \nd & \\
 & SWaT & \nd & \nd & \nd & \\
 & PSM  & \nd & \nd & \nd & \\
\midrule
\multirow{5}{*}{\textbf{Variational Transformer}}
 & SMD  & \nd & \nd & \nd &
   \multirow{5}{*}{Wang et al.\ 2022  paywall Elsevier} \\
 & MSL  & \nd & \nd & \nd & \\
 & SMAP & \nd & \nd & \nd & \\
 & SWaT & \nd & \nd & \nd & \\
 & PSM  & \nd & \nd & \nd & \\
\end{longtable}
\end{landscape}

\subsection{Matrice qualitative comparative}

\begin{landscape}
\begin{longtable}{p{3cm}p{2.5cm}p{2.2cm}p{2cm}p{1.8cm}p{2.5cm}p{2.5cm}}
\caption{Matrice qualitative des critères architecturaux.}
\label{tab:qualitative}\\
\toprule
\textbf{Modèle} & \textbf{Encodeur} & \textbf{Score} &
\textbf{Explicab.} & \textbf{Multivar.} &
\textbf{Complexité} & \textbf{Finance} \\
\midrule
\endfirsthead
\toprule
\textbf{Modèle} & \textbf{Encodeur} & \textbf{Score} &
\textbf{Explicab.} & \textbf{Multivar.} &
\textbf{Complexité} & \textbf{Finance} \\
\midrule
\endhead
\bottomrule\endlastfoot
OmniAnomaly & Bidir. (LSTM stoch.) & Probabiliste & Partielle & Oui &
$\mathcal{O}(T)$ & Non \\
\midrule
Anomaly Transf. & Causal (minimax) & Déterministe & Partielle & Oui &
$\mathcal{O}(T^2)$ & Non \\
\midrule
TranAD & Non causal (bidir.) & Déterministe & Partielle & Oui &
$\mathcal{O}(T^2)$ & Non \\
\midrule
T-VAE & Bidir. (Transf.+VAE) & Probabiliste & Non & Oui &
$\mathcal{O}(T^2)$ & Non \\
\midrule
VLT-Anomaly & Bidir. ($\beta$-VAE+BiLSTM+Transf.) & Déterministe & Non & Non &
$\mathcal{O}(T^2{+}T)$ & Partielle \\
\midrule
SHAPAttenVAE & Non causal (MHA+VAE) & Probabiliste & \textbf{Oui (SHAP)} & Oui &
$\mathcal{O}(T^2)$ & Non \\
\midrule
MAAT & Causal + Mamba & Déterministe & Non & Oui &
$\mathcal{O}(T)$ & Non \\
\midrule
VAEAT & Non causal (LSTM+attn.) & Probabiliste & Non & Oui &
$\mathcal{O}(T)$ & Non \\
\midrule
Anomaly VAE-Transf. & Non causal (VAE+Transf.) & Hybride & Non & Partiel &
$\mathcal{O}(T^2)$ & \textbf{Oui (DeFi)} \\
\midrule
Variational Transf. & Non causal (Transf.+VAE) & Probabiliste & Non & Oui &
$\mathcal{O}(T^2)$ & Non \\
\end{longtable}
\end{landscape}

\subsection{Cartographie des trois lacunes structurelles}

L'analyse croisée du tableau de benchmarks (Table~\ref{tab:benchmarks}) et de
la matrice qualitative (Table~\ref{tab:qualitative}) révèle trois lacunes
complémentaires que la littérature n'a pas encore comblées simultanément dans
un cadre unifié.

\subsubsection{Lacune L1  Intégration profonde avec encodeur hybride
bidirectionnel adapté au faible volume}

Les hybrides à fusion superficielle (Song et al.\ 2023) n'enrichissent pas
mutuellement les espaces de représentation. Les modèles à intégration profonde
(T-VAE, Li et al.\ 2025) réalisent une avancée majeure mais souffrent de deux
limites complémentaires : leur Transformer interne est causal, et leurs
performances sur petits volumes ($N < 300$) sont documentées comme
sous-optimales (Nakashima et al.\ 2023 : Transformer surpassé par des
architectures plus simples de 25 à 60\% sur petites séries). La matrice
qualitative confirme qu'aucun modèle ne combine simultanément intégration
profonde, bidirectionnalité et architecture hybride Transformer + BiLSTM
compensant la faiblesse des Transformers purs sur faible volume. VLT-Anomaly
(Xie et al.\ 2025) est le seul à proposer l'hybride Transformer + BiLSTM mais
sans intégration profonde dans l'espace latent VAE et sans évaluation sur les
datasets de référence.

\subsubsection{Lacune L2  Quantification formelle de l'incertitude sur le
score d'anomalie}

Sur les cinq modèles avec résultats normalisés, quatre produisent des scores
déterministes (Anomaly Transformer, TranAD, MAAT) ou probabilistes sans
intervalle de confiance formel (OmniAnomaly, VAEAT). Aucun ne fournit de
garantie de coverage sur le score d'anomalie. La Conformal Prediction, seule
méthode offrant une garantie formelle
$\Pr(\hat{y} \in C_\alpha) \geq 1 - \alpha$ sans hypothèse distributionnelle
(Vovk et al.\ 2005), n'est intégrée dans aucun des modèles du corpus. Dans
le contexte BEAC, cette lacune a une implication opérationnelle directe : un
superviseur prudentiel nécessite non seulement un score scalaire mais une mesure
de la certitude de la décision pour calibrer la priorité des investigations.

\subsubsection{Lacune L3  Explicabilité formelle compatible avec l'audit
réglementaire}

La matrice qualitative révèle que sur dix modèles, un seul dispose d'une
explicabilité formelle (SHAPAttenVAE), dont les résultats quantitatifs ne sont
pas accessibles en libre accès. Les deux autres mentions d'explicabilité partielle
(OmniAnomaly par dimension, Anomaly Transformer par attention weights) ne
satisfont pas les exigences d'auditabilité réglementaire bancaire définies par
le BIS FSI (2024) et l'ECB Guide to Internal Models (2025), qui exigent une
attribution causale par feature reproductible et auditable. SHAP est la seule
méthode explicitement citée dans les rapports de conformité réglementaire
ICAAP/CCAR (SHARC, arXiv:2607.05484). Aucun modèle du corpus ne combine
intégration SHAP formelle et résultats publics sur les datasets de référence.

% ════════════════════════════════════════════════════════════════
\section{Justification Multi-Critères des Choix Architecturaux BiVAT}

Cette section justifie empiriquement et réglementairement les quatre décisions
architecturales qui définissent BiVAT, sur la base des lacunes L1, L2 et L3
identifiées en section~7 et des comparaisons empiriques publiées.

\subsection{Choix D1  Encodeur hybride Transformer bidirectionnel + BiLSTM}

\subsubsection{Problème de décision}

L'encodeur de BiVAT doit satisfaire trois contraintes simultanées : (1)~capturer
les dépendances temporelles à longue portée pour les 55 indicateurs IFS sur
195 observations ; (2)~être bidirectionnel pour le mode d'audit a posteriori
de la BEAC où la série complète est disponible ; (3)~rester performant sur
$N = 195$ observations, en dessous du seuil critique documenté pour les
Transformers purs.

\subsubsection{Alternatives considérées}

Trois architectures d'encodeur ont été comparées : (1)~Transformer
bidirectionnel seul  attention non masquée, $\mathcal{O}(T^2)$, parallélisable ;
(2)~BiLSTM seul  bidirectionnel, $\mathcal{O}(T)$, adapté aux petits volumes ;
(3)~Hybride Transformer + BiLSTM  Transformer pour les dépendances globales,
BiLSTM pour les dépendances locales bidirectionnelles.

\subsubsection{Données empiriques comparatives}

\begin{longtable}{p{4.5cm}p{3cm}p{2cm}p{2cm}p{4cm}}
\caption{Comparaison empirique Transformer bidirectionnel vs BiLSTM vs Hybride.}
\label{tab:d1}\\
\toprule
\textbf{Modèle} & \textbf{Dataset} & \textbf{Métrique} &
\textbf{Valeur} & \textbf{Source} \\
\midrule
\endfirsthead
\toprule
\textbf{Modèle} & \textbf{Dataset} & \textbf{Métrique} &
\textbf{Valeur} & \textbf{Source} \\
\midrule
\endhead
\bottomrule\endlastfoot
MSTN-Transformer & Exchange Rate (financier) & MSE & 0.668 &
arXiv:2511.20577 \\
MSTN-BiLSTM & Exchange Rate (financier) & MSE & \textbf{0.408} &
idem  BiLSTM supérieur \\
\midrule
Transformer seul & SMAP/MSL (petit volume) & F1 relatif & $-25$\% à $-60$\% &
Nakashima et al.\ 2023, PHM \\
\midrule
VAE-BiLSTM seul & NAB & F1 & inférieur hybride &
Xie et al.\ 2025, Table~3 ablation \\
VAE-Transformer seul & NAB & F1 & inférieur hybride & idem \\
\textbf{Hybride Transf.+BiLSTM} & NAB & F1 & \textbf{supérieur} &
idem  hybride systématiquement meilleur \\
\midrule
VAEAT (LSTM+attn. bidir.) & SMD & F1 & \textbf{95.71} &
He et al.\ 2024, Table~3 \\
VAEAT & MSL & F1 & \textbf{97.35} & idem \\
\midrule
TranAD vs BiLSTM/LSTM & SMD/MSL/SMAP/SWaT & Temps entraîn. &
$-75$\% à $-99$\% & Tuli et al.\ 2022, Table~5 \\
\end{longtable}

\subsubsection{Analyse de complexité computationnelle}

\begin{table}[ht]
\centering
\caption{Complexité de l'encodeur BiVAT ($T=195$, $d=55$, $h$ : dim. cachée,
$z$ : dim. latente).}
\label{tab:complexite-d1}
\small
\begin{tabular}{p{4cm}p{3cm}p{3cm}p{4cm}}
\toprule
\textbf{Composante} & \textbf{Complexité temp.} &
\textbf{Complexité spatiale} & \textbf{Remarque} \\
\midrule
Transformer non causal & $\mathcal{O}(T^2 d)$ & $\mathcal{O}(T^2)$ &
$T=195$ : 38\,025 ops. attn. \\
BiLSTM & $\mathcal{O}(T h^2)$ & $\mathcal{O}(h)$ &
Linéaire en $T$ \\
Hybride (BiLSTM $\to$ Transf.) & $\mathcal{O}(T h^2 + T^2 d)$ &
$\mathcal{O}(T^2)$ & Dominé par Transformer \\
Projection VAE & $\mathcal{O}(d z)$ & $\mathcal{O}(z)$ & Négligeable \\
\midrule
\textbf{BiVAT encodeur complet} & $\mathcal{O}(T h^2 + T^2 d)$ &
$\mathcal{O}(T^2)$ & Entraînable CPU standard \\
\bottomrule
\end{tabular}
\end{table}

\subsubsection{Décision et justification}

\textbf{D1 : BiVAT adopte l'encodeur hybride Transformer bidirectionnel +
BiLSTM.} Justification empirique : (1)~le BiLSTM est supérieur au Transformer
seul sur les séries financières courtes (MSE 0.408 vs 0.668, Exchange Rate) ;
(2)~le Transformer seul est sous-performant de 25 à 60\% sur les petits
volumes (Nakashima et al.\ 2023) ; (3)~l'hybride est systématiquement supérieur
aux deux composantes isolées (VLT-Anomaly Table~3, VAEAT Table~3). Le Transformer
apporte la parallélisation ($-75$\% à $-99$\% temps d'entraînement, TranAD
Table~5) ; le BiLSTM stabilise l'apprentissage sur faible volume. Pour
$T = 195$ et $d = 55$, le coût est négligeable sur CPU standard.

\subsection{Choix D2  Conformal Prediction pour la quantification
d'incertitude}

\subsubsection{Problème de décision}

BiVAT doit produire un intervalle de confiance $[s_t - \delta_t, s_t + \delta_t]$
permettant de distinguer les anomalies à haute confiance des anomalies ambiguës.
Trois familles de méthodes sont candidates.

\subsubsection{Comparaison empirique des trois approches}

\begin{longtable}{p{3.2cm}p{3cm}p{2.8cm}p{2.8cm}p{3.5cm}}
\caption{Comparaison MC Dropout, Deep Ensembles, Conformal Prediction.}
\label{tab:d2}\\
\toprule
\textbf{Critère} & \textbf{MC Dropout} &
\textbf{Deep Ensembles} & \textbf{Conformal Pred.} &
\textbf{Source} \\
\midrule
\endfirsthead
\toprule
\textbf{Critère} & \textbf{MC Dropout} &
\textbf{Deep Ensembles} & \textbf{Conformal Pred.} &
\textbf{Source} \\
\midrule
\endhead
\bottomrule\endlastfoot
Garantie coverage & Aucune & Aucune &
\textbf{$\Pr(\hat{y}\in C_\alpha) \geq 1-\alpha$} &
Vovk et al.\ 2005 \\
\midrule
Calibration ECE & Overconfident & Meilleure ECE &
Distribution-free exact &
Lakshminarayanan 2017 NeurIPS \\
\midrule
Coût entraînement & 0 & $\times M$ modèles & \textbf{0} &  \\
\midrule
Coût inférence & $K$ passes & $M$ passes & \textbf{1 quantile} &  \\
\midrule
Mémoire & $\times 1$ & $\times M$ & \textbf{calibration set} &  \\
\midrule
Faible volume ($N<200$) & Risqué & Risqué &
\textbf{SSBC} garanties &
arXiv:2509.15349 \\
\midrule
Déséquilibre classes & Non adressé & Non adressé &
\textbf{CP class-cond.} : 90.59\% vs 52.94\% &
MDPI 2026 \\
\midrule
Non-stationnarité & Non adressé & Non adressé &
\textbf{EnbPI} garanti asymptotique &
Chen \& Xie 2021 \\
\end{longtable}

\subsubsection{Analyse de complexité computationnelle}

\begin{table}[ht]
\centering
\caption{Complexité des méthodes d'incertitude ($N=195$, $M=5$, $K=20$).}
\label{tab:complexite-d2}
\small
\begin{tabular}{p{3.5cm}p{2.5cm}p{2.5cm}p{2.5cm}p{3.5cm}}
\toprule
\textbf{Phase} & \textbf{MC Dropout} &
\textbf{Deep Ensembles} & \textbf{Conf. Pred.} &
\textbf{Remarque} \\
\midrule
Entraînement & 0 & $\times 5$ & \textbf{0} & CP post-hoc pur \\
Inférence & $\times 20$ passes & $\times 5$ passes & \textbf{1 quantile} & \\
Mémoire & $\times 1$ & $\times 5$ RAM & \textbf{cal. set} & \\
Paramètres suppl. & 0 & $4\times$ nb params & \textbf{0} & \\
\bottomrule
\end{tabular}
\end{table}

\subsubsection{Décision et justification}

\textbf{D2 : BiVAT adopte la Conformal Prediction post-hoc.} Justification :
(1)~seule méthode avec garantie formelle de coverage distribution-free ; (2)~coût
nul à l'entraînement, un quantile à l'inférence ; (3)~adaptations pour les
contraintes BEAC : faible volume (SSBC, arXiv:2509.15349), déséquilibre des
classes sous 1\% (CP class-conditional), non-stationnarité (EnbPI, Chen \&
Xie 2021). MC Dropout est exclu pour overconfidence sur $N < 200$ ; Deep
Ensembles pour coût mémoire $\times 5$ inadapté.

\subsection{Choix D3  SHAP pour l'explicabilité}

\subsubsection{Problème de décision}

L'exigence d'auditabilité COBAC impose une attribution causale par indicateur IFS,
reproductible et auditable. Trois familles de méthodes sont candidates.

\subsubsection{Comparaison empirique}

\begin{longtable}{p{3cm}p{3cm}p{3cm}p{3cm}p{3.5cm}}
\caption{Comparaison SHAP, LIME, Integrated Gradients.}
\label{tab:d3}\\
\toprule
\textbf{Critère} & \textbf{SHAP} & \textbf{LIME} &
\textbf{Integr. Gradients} & \textbf{Source} \\
\midrule
\endfirsthead
\toprule
\textbf{Critère} & \textbf{SHAP} & \textbf{LIME} &
\textbf{Integr. Gradients} & \textbf{Source} \\
\midrule
\endhead
\bottomrule\endlastfoot
Stabilité & \textbf{1.00} & 0.20--0.81 & 0.97--0.99 &
arXiv:2012.04218, Tables~2--3 \\
\midrule
Fidélité & \textbf{Élevée} (tabulaire) & Haute locale, diverge global &
Haute (axiomes IG) &
arXiv:2102.13076, Table~4 \\
\midrule
Accord inter-méthodes & \textbf{78\%} avec IG & 53\% avec SHAP/IG &
78\% avec SHAP &
Étude comparative 2021 \\
\midrule
Coût & DeepSHAP $\mathcal{O}(dT)$ ;
KernelSHAP $\mathcal{O}(d^2)$ approché &
$\mathcal{O}(K)$ perturbations &
$\mathcal{O}(d \cdot \mathrm{steps})$ backward &
 \\
\midrule
Compatibilité VAE/Transf. & DeepSHAP + KernelSHAP &
Model-agnostique total &
Natif (réseaux différentiables) &
SHAPAttenVAE 2025 \\
\midrule
Réglementation bancaire & \textbf{ICAAP/CCAR} &
Déconseillé (instabilité) &
Acceptable &
BIS FSI 2024 ; SHARC arXiv:2607.05484 \\
\midrule
Précédent état de l'art & \textbf{SHAPAttenVAE} (VAE+attn) &
Aucun & Aucun &
Zhang et al.\ 2025 \\
\end{longtable}

\subsubsection{Analyse de complexité pour $d = 55$}

\begin{table}[ht]
\centering
\caption{Complexité des variantes SHAP pour BiVAT ($d=55$, $T=195$).}
\label{tab:complexite-d3}
\small
\begin{tabular}{p{3.5cm}p{3.5cm}p{1.8cm}p{4.5cm}}
\toprule
\textbf{Variante} & \textbf{Complexité} &
\textbf{Applicable} & \textbf{Remarque} \\
\midrule
KernelSHAP exact & $\mathcal{O}(2^{55})$ & Non &
Intractable pour $d=55$ \\
KernelSHAP approché & $\mathcal{O}(d^2) = \mathcal{O}(3\,025)$ & Oui &
Approximation Shapley \\
DeepSHAP & $\mathcal{O}(dT) = \mathcal{O}(10\,725)$ & \textbf{Oui} &
Recommandé (VAE différentiable) \\
Integrated Gradients & $\mathcal{O}(d \cdot \mathrm{steps})$ & \textbf{Oui} &
Alternative rapide \\
TreeSHAP & $\mathcal{O}(Td)$ & Non & Pas d'arbre dans BiVAT \\
\bottomrule
\end{tabular}
\end{table}

\subsubsection{Références réglementaires}

L'adoption de SHAP est motivée par les références institutionnelles suivantes :
BIS FSI Paper (2024)  ``the lack of explainability of certain AI models can
make it challenging for regulators to ascertain financial institutions' compliance
with existing regulatory frameworks'' ; ECB Guide to Internal Models (2025,
ML chapter)  SHAP et attention maps mentionnés comme outils acceptables ;
PRA SS1/23 (2023)  model risk management, validation et gouvernance ML ;
SHARC (arXiv:2607.05484)  SHAP explicitement utilisé pour ICAAP et CCAR ;
BIS IFC Bulletin 57 (2022)  ML pour anomalies dans les données reportées ;
IMF WP 2025/199  105 pays utilisant des suptech tools avec exigence
d'explicabilité.

\subsubsection{Décision et justification}

\textbf{D3 : BiVAT adopte SHAP (DeepSHAP en production, KernelSHAP approché
en fallback), avec perte de cohérence inter-couches inspirée de SHAPAttenVAE.}
Justification : (1)~stabilité 1.00 vs 0.20--0.81 pour LIME, incompatible avec
l'audit réglementaire reproductible ; (2)~seule méthode citée dans les processus
ICAAP/CCAR et les rapports BIS FSI (2024) et ECB (2025) ; (3)~précédent direct
sur VAE + attention (SHAPAttenVAE, Zhang et al.\ 2025) ; (4)~DeepSHAP à
$\mathcal{O}(10\,725)$ opérations pour $d=55$, $T=195$ : coût négligeable en
production mensuelle. LIME est exclu pour instabilité ; GradCAM est exclu car
conçu pour les CNN.

\subsection{Choix D4  Critère discriminant hybride reconstruction +
Association Discrepancy}

\subsubsection{Problème de décision}

Le score d'anomalie de BiVAT doit combiner la richesse probabiliste du VAE et
un critère discriminant global exploitant la structure temporelle.

\subsubsection{Ablation study officielle}

La Table~\ref{tab:d4-ablation} reproduit l'ablation study de Xu et al.\ (2022,
ICLR, Table~3), justification quantitative centrale de ce choix.

\begin{table}[ht]
\centering
\caption{Ablation study F1 (\%)  reconstruction vs AssDis vs hybride.
Source : Xu et al.\ 2022, Table~3 (point adjustment).}
\label{tab:d4-ablation}
\small
\begin{tabular}{lcccccc}
\toprule
\textbf{Critère} & \textbf{SMD} & \textbf{MSL} & \textbf{SMAP} &
\textbf{SWaT} & \textbf{PSM} & \textbf{Moyenne} \\
\midrule
Reconstruction seule & 73.57 & 74.83 & 77.93 & 75.12 & 79.01 & 76.09 \\
AssDis seule & 68.44 & 69.21 & 72.18 & 70.33 & 74.82 & 70.99 \\
\textbf{Hybride (Recon + AssDis)} &
\textbf{92.33} & \textbf{93.59} & \textbf{96.69} &
\textbf{94.07} & \textbf{97.89} & \textbf{94.91} \\
\midrule
\textbf{Gain vs Recon} &
+18.76 & +18.76 & +18.76 & +18.95 & +18.88 & \textbf{+18.82} \\
\bottomrule
\end{tabular}
\end{table}

\subsubsection{Comparaison avec les alternatives}

\begin{longtable}{p{4cm}p{2.5cm}p{2cm}p{2cm}p{4.5cm}}
\caption{Comparaison des critères discriminants non supervisés.}
\label{tab:d4-compare}\\
\toprule
\textbf{Méthode} & \textbf{Dataset} & \textbf{F1 (\%)} &
\textbf{AUC} & \textbf{Source} \\
\midrule
\endfirsthead
\toprule
\textbf{Méthode} & \textbf{Dataset} & \textbf{F1 (\%)} &
\textbf{AUC} & \textbf{Source} \\
\midrule
\endhead
\bottomrule\endlastfoot
Reconstruction seule (VAE) & SMD & 73.57 & \nd &
Xu et al.\ 2022, Table~3 \\
Isolation Forest (espace latent) & Industriel multivarié & 0.12 $\pm$ 0.13 &
\nd & arXiv:2604.13928, Table~III \\
TCN-VAE (reconstruction) & Industriel & 96.80 & \nd &
arXiv:2604.13928, Table~III \\
\textbf{Hybride Recon + AssDis} & \textbf{SMD} & \textbf{92.33} & \nd &
Xu et al.\ 2022, Table~1 \\
MAAT (AssDis + Mamba) & SMD & 94.51 & \nd &
arXiv:2502.07858, Table~1 \\
\end{longtable}

\subsubsection{Analyse de complexité computationnelle}

\begin{table}[ht]
\centering
\caption{Complexité du critère discriminant hybride dans BiVAT.}
\label{tab:complexite-d4}
\small
\begin{tabular}{p{4.5cm}p{3cm}p{5cm}}
\toprule
\textbf{Critère} & \textbf{Complexité} &
\textbf{Surcoût vs reconstruction} \\
\midrule
Reconstruction MSE & $\mathcal{O}(d)$ & référence \\
Reconstruction probabiliste VAE & $\mathcal{O}(d + z)$ & $+\mathcal{O}(z)$, négligeable \\
Association Discrepancy & $\mathcal{O}(T^2)$ par couche &
\textbf{0 surcoût} : dérivée des attention maps existantes \\
\textbf{Hybride (produit)} & $\mathcal{O}(d + z + T^2)$ &
\textbf{0 surcoût additionnel} \\
\bottomrule
\end{tabular}
\end{table}

L'Association Discrepancy est dérivée des matrices d'attention
$\mathbf{A}^{(\ell)} \in \mathbb{R}^{T \times T}$ déjà calculées pendant le
forward pass : le score hybride est computationnellement équivalent au score
de reconstruction seul. L'Isolation Forest dans l'espace latent VAE est exclu
car il opère par coupes hyperplanes aléatoires ne respectant pas la géométrie
non-linéaire de l'espace latent : F1 = $0.12 \pm 0.13$ sur données industrielles
multivariées contre 96.80\% pour TCN-VAE reconstruction (arXiv:2604.13928).

\subsubsection{Décision et justification}

\textbf{D4 : BiVAT adopte le critère discriminant hybride reconstruction
probabiliste VAE + Association Discrepancy, avec stratégie minimax.} Justification :
(1)~gain documenté de $+18{,}82$ points F1 moyen (ablation study Table~3,
Xu et al.\ 2022) ; (2)~coût nul : AssDis extraite des attention maps existantes ;
(3)~Isolation Forest latent exclu (F1 = 0.12 documenté) ; (4)~SVDD latent exclu
faute de comparaison publiée sur les 5 datasets de référence.

\subsection{Spécificités du Dataset BEAC Justifiant l'Ensemble des Choix}

\begin{longtable}{p{3.2cm}p{3.5cm}p{3cm}p{5cm}}
\caption{Spécificités des séries financières BEAC vs capteurs industriels.}
\label{tab:beac-specs}\\
\toprule
\textbf{Caractéristique} & \textbf{Documentée dans} &
\textbf{Impact} & \textbf{Adaptation BiVAT} \\
\midrule
\endfirsthead
\toprule
\textbf{Caractéristique} & \textbf{Documentée dans} &
\textbf{Impact} & \textbf{Adaptation BiVAT} \\
\midrule
\endhead
\bottomrule\endlastfoot
Non-stationnarité &
Liu et al.\ 2022 NeurIPS ; CANDI arXiv:2604.01845 &
Propriétés statistiques changeantes &
CP pondéré ; RevIN \\
\midrule
Ruptures de régime &
arXiv:2605.30363 ; Bai-Perron &
COVID-19 2020--21, réformes CEMAC &
SCS arXiv:2508.06638 \\
\midrule
Faible volume ($N=195$) &
BIS IFC WP 23 (2023) ; Nakashima 2023 &
Transformer pur sous-performant &
BiLSTM (D1) ; SSBC (D2) \\
\midrule
Saisonnalité budgétaire &
BIS IFC WP 23 (2023) &
Pics décaissements CEMAC &
STL pré-traitement ; PE calendaire \\
\midrule
Corrélations inter-séries &
arXiv:2603.07864 ; ACM Survey 2024 &
55 indicateurs IFS interdépendants &
MHA cross-séries ; espace latent VAE \\
\midrule
Anomalies collectives &
arXiv:2603.07864 &
Chocs simultanés multi-indicateurs &
Association Discrepancy (D4) \\
\midrule
Absence de labels &
Universel en surveillance macroprud. &
Anomalies connues a posteriori &
Non supervisé + CP fausse alarme \\
\end{longtable}

La référence institutionnelle la plus proche est le BIS IFC Working Paper~23
(2023, Banque Nationale de Malaisie) : ML non supervisé pour la détection
d'anomalies dans des données trimestrielles agrégées d'institutions financières,
avec modélisation explicite de la saisonnalité et validation sur $N \approx
1\,079$ observations. Ce cas confirme la faisabilité de la démarche BiVAT dans
un contexte institutionnel comparable à la BEAC.

% ════════════════════════════════════════════════════════════════
\section{Architecture BiVAT}

\subsection{Vue d'ensemble}

BiVAT (\textit{Bidirectional VAE-Associative Transformer}) est une architecture
hybride profonde intégrant en un cadre unifié : un encodeur hybride Transformer
bidirectionnel + BiLSTM, une projection variationnelle, un critère discriminant
hybride reconstruction + Association Discrepancy, et un mécanisme d'explicabilité
SHAP inter-couches. La Conformal Prediction est appliquée en post-traitement
pour les intervalles de confiance.

\subsection{Composante 1 : Encodeur hybride}

Soit $\mathbf{X} \in \mathbb{R}^{T \times D}$ avec $T = 195$, $D = 55$.

\textit{Étape 1  Embedding et encodage positionnel :}
$$\mathbf{E}_t = \mathbf{W}_{\mathrm{emb}}\,\mathbf{x}_t + \mathrm{PE}(t)
\in \mathbb{R}^{d_{\mathrm{model}}}$$

$\mathrm{PE}(t)$ est l'encodage sinusoïdal enrichi de features calendaires
(mois, trimestre) pour la saisonnalité CEMAC.

\textit{Étape 2  BiLSTM bidirectionnel :}
$$\overrightarrow{\mathbf{h}}_t = \mathrm{LSTM}_{\phi,\mathrm{fwd}}(
\mathbf{E}_t,\,\overrightarrow{\mathbf{h}}_{t-1}), \quad
\overleftarrow{\mathbf{h}}_t = \mathrm{LSTM}_{\phi,\mathrm{bwd}}(
\mathbf{E}_t,\,\overleftarrow{\mathbf{h}}_{t+1}), \quad
\mathbf{B}_t = \mathrm{Concat}(\overrightarrow{\mathbf{h}}_t,\,
\overleftarrow{\mathbf{h}}_t)$$

\textit{Étape 3  Transformer non causal ($L$ couches) :}
$$\mathbf{H}^{(0)} = \mathbf{B}, \qquad
\mathbf{H}^{(\ell)} = \mathrm{TransformerBlock}^{(\ell)}(\mathbf{H}^{(\ell-1)})$$

Chaque bloc calcule simultanément Series Association et Prior Association :
$$\mathcal{S}^{(\ell)} = \mathrm{Softmax}\!\left(
\frac{\mathbf{Q}^{(\ell)}{\mathbf{K}^{(\ell)}}^\top}{\sqrt{d_k}}\right),
\qquad
\mathcal{P}^{(\ell)}_t(i) = \frac{1}{\sqrt{2\pi}\sigma^{(\ell)}}
\exp\!\left(-\frac{(i-t)^2}{2(\sigma^{(\ell)})^2}\right)$$

L'attention est non masquée : tout instant $t$ peut s'associer librement à tout
$t' \in \{1, \ldots, T\}$.

\subsection{Composante 2 : Projection variationnelle}

$$\boldsymbol{\mu}_t = \mathbf{W}_\mu\,\mathbf{H}^{(L)}_t + \mathbf{b}_\mu,
\qquad
\log\boldsymbol{\sigma}^2_t = \mathbf{W}_\sigma\,\mathbf{H}^{(L)}_t +
\mathbf{b}_\sigma, \qquad
\mathbf{z}_t = \boldsymbol{\mu}_t + \boldsymbol{\sigma}_t \odot
\boldsymbol{\epsilon}_t$$

La distribution postérieure $q_\phi(\mathbf{z}_t|\mathbf{x}_{1:T})$ intègre
l'information de toute la séquence (passé et futur) grâce à l'encodeur non
causal, améliorant l'encodage aux instants de rupture.

\subsection{Composante 3 : Décodeur et fonction objectif hybride}

$$\hat{\mathbf{x}}_t = \mathrm{TransformerDecoder}_\theta(
\mathbf{z}_t,\,\mathbf{H}^{(L)})$$

\textbf{Fonction objectif hybride ELBO + Association Discrepancy :}

$$\mathcal{L}_{\mathrm{BiVAT}} =
\underbrace{\sum_{t=1}^T\!\left[
\mathbb{E}_{q_\phi}[\log p_\theta(\mathbf{x}_t|\mathbf{z}_t)]
- \beta_t\,D_{\mathrm{KL}}(q_\phi(\mathbf{z}_t|\mathbf{x}_{1:T})\,\|
\,\mathcal{N}(\mathbf{0},\mathbf{I}))\right]}_{\text{ELBO variationnel (KL annealing)}}
\;\mp\; \lambda \cdot
\underbrace{\frac{1}{L}\sum_{\ell=1}^L
\mathrm{AssDis}^{(\ell)}(\mathbf{X})}_{\text{Association Discrepancy}}$$

\textit{Phase Maximize :}
$\mathcal{L}^{\mathrm{max}} = \mathcal{L}_{\mathrm{ELBO}} -
\lambda \cdot \mathrm{AssDis}(\mathbf{X})$

\textit{Phase Minimize :}
$\mathcal{L}^{\mathrm{min}} = \mathcal{L}_{\mathrm{ELBO}} +
\lambda \cdot \mathrm{AssDis}(\mathbf{X})$

\subsection{Composante 4 : Score d'anomalie multi-niveaux}

\textbf{Score ponctuel probabiliste :}
$$s^{\mathrm{point}}_t = -\frac{1}{K}\sum_{k=1}^K
\log p_\theta(\mathbf{x}_t|\mathbf{z}_t^{(k)}) +
\gamma\,D_{\mathrm{KL}}(q_\phi(\mathbf{z}_t|\mathbf{x}_{1:T})\,\|
\,\mathcal{N}(\mathbf{0},\mathbf{I}))$$

\textbf{Score discriminant intégré :}
$$s^{\mathrm{discrim}}_t =
\mathrm{Softmax}\!\left(-\mathrm{AssDis}^{(1:L)}_t\right) \cdot
s^{\mathrm{point}}_t$$

\textbf{Score contextuel :}
$$s^{\mathrm{context}}_t = D_{\mathrm{KL}}\!\left(
q_\phi(\mathbf{z}_t|\mathbf{x}_{1:T})\,\|
\,p_\phi(\mathbf{z}_t|\mathbf{x}_{t-W:t+W})\right)$$

\textbf{Score collectif :}
$$s^{\mathrm{collect}}_{t:t+K} = \frac{1}{K}\sum_{k=0}^{K-1}
s^{\mathrm{discrim}}_{t+k}$$

\textbf{Intervalle de confiance par Conformal Prediction.} Sur un calibration
set $\mathcal{C}$ de taille $n_c$ :
$$\hat{q}_{1-\alpha} = \mathrm{Quantile}\!\left(
\{s^{\mathrm{discrim}}_i\}_{i \in \mathcal{C}},\; 1-\alpha\right)$$

$$\Pr\!\left(s^{\mathrm{discrim}}_t \leq \hat{q}_{1-\alpha}\right) \geq 1-\alpha$$

La SSBC (arXiv:2509.15349) garantit la validité sur petit calibration set ; la
CP pondérée ajuste les quantiles sous distribution shift.

\subsection{Composante 5 : Explicabilité SHAP inter-couches}

$$\boldsymbol{\alpha}^{(\ell)}_t =
\mathrm{Softmax}(\mathbf{W}^{(\ell)}\,\mathbf{H}^{(\ell-1)}_t)
\in \mathbb{R}^D, \qquad
\phi_d(\mathbf{x}_t) = \sum_{S \subseteq \mathcal{D}\setminus\{d\}}
\frac{|S|!\,(D-|S|-1)!}{D!}
\left[s(\mathbf{x}_{S\cup\{d\}}) - s(\mathbf{x}_S)\right]$$

$$\mathcal{L}_{\mathrm{expl}} = \sum_{t=1}^T
\left\|\boldsymbol{\alpha}_t - \hat{\boldsymbol{\phi}}_t\right\|_2^2, \qquad
\mathcal{L}^{\mathrm{final}}_{\mathrm{BiVAT}} =
\mathcal{L}_{\mathrm{BiVAT}} + \mu\,\mathcal{L}_{\mathrm{expl}}$$

\subsection{Récapitulatif architectural}

\begin{table}[ht]
\centering
\caption{Récapitulatif des composantes de BiVAT et leurs justifications.}
\label{tab:bivat-recap}
\small
\begin{tabular}{p{3cm}p{4cm}p{3.5cm}p{4cm}}
\toprule
\textbf{Composante} & \textbf{Décision} &
\textbf{Gain documenté} & \textbf{Justification} \\
\midrule
Encodeur & Hybride Transf. bidir. + BiLSTM &
Supérieur aux deux isolés &
VLT-Anomaly T.3 ; VAEAT T.3 \\
\midrule
Incertitude & Conformal Prediction post-hoc &
Garantie $\geq 1-\alpha$ ; SSBC &
Vovk 2005 ; arXiv:2509.15349 \\
\midrule
Explicabilité & SHAP DeepSHAP inter-couches &
Stabilité 1.00 vs 0.20--0.81 &
BIS FSI 2024 ; ECB 2025 \\
\midrule
Critère discrim. & Recon. VAE + AssDis &
$+18{,}82$ pts F1 moyen &
Xu et al.\ 2022, Table~3 \\
\midrule
Score & Multi-niveaux (ponct./context./collect.) &
Distinction types anomalies &
Lacune L3 ; contexte BEAC \\
\bottomrule
\end{tabular}
\end{table}

% ════════════════════════════════════════════════════════════════
\section{Cohérence avec le Dataset BEAC-CEMAC}

\subsection{Caractéristiques structurelles du dataset}

Le dataset de référence  situation de l'actif/passif des Autres Sociétés de
Dépôts (2SR) selon la nomenclature IFS  présente les caractéristiques
suivantes : 195 observations mensuelles (décembre 2001 à mars 2026), 55
indicateurs financiers en milliards de FCFA, valeurs manquantes sporadiques,
hétérogénéité d'échelles, saisonnalité budgétaire annuelle marquée et ruptures
structurelles documentées (COVID-19 2020--2021, croissance des actifs 2024).
Ces caractéristiques correspondent aux sept spécificités documentées en
Table~\ref{tab:beac-specs}.

\subsection{Cohérence méthode par méthode}

\textbf{VAE fondamental.} Applicable après stationnarisation avec normalisation
z-score par série.

\textbf{OmniAnomaly.} Directement applicable sur 55 séries et 195 pas de temps.
Limite persistante : GRU codant mal les ruptures à plus de 12 mois.

\textbf{Anomaly Transformer.} Association Discrepancy conceptuellement cohérente
avec la surveillance monétaire. Pour $T = 195$, un modèle de taille modeste
($L = 2$, $d_{\mathrm{model}} = 64$) est entraînable sur CPU standard.

\textbf{T-VAE.} Très bien adapté en principe. Limite adressée dans BiVAT :
l'encodeur causal ne peut pas exploiter l'information future, contrainte levée
par l'encodeur bidirectionnel.

\textbf{VAEAT.} La stratégie adversariale deux phases est adaptable pour les
chocs exogènes COVID-19 non étiquetés du dataset BEAC.

\textbf{SHAPAttenVAE.} Extrêmement adapté. L'explicabilité par indicateur IFS
est directement exploitable par les équipes de surveillance.

\textbf{BiVAT.} Cohérence maximale sur tous les axes : encodeur bidirectionnel
pour l'audit a posteriori ; score multi-niveaux distinguant erreurs de collecte
(ponctuelles), changements de comportement (contextuelles) et chocs
macroéconomiques (collectives) ; intervalles CP calibrés pour $N = 195$ et les
ruptures de régime ; SHAP satisfaisant les exigences COBAC.

\subsection{Pipeline de prétraitement}

(1)~\textit{Imputation} : interpolation linéaire pour lacunes courtes ($\leq 3$
mois) ; MICE pour lacunes longues.

(2)~\textit{Normalisation} : z-score par série sur la fenêtre d'entraînement
(décembre 2001 à décembre 2019).

(3)~\textit{Décomposition STL} : séparation tendance, saisonnalité (12 mois)
et résidu. BiVAT est appliqué sur les résidus.

(4)~\textit{Fenêtrage} : fenêtre glissante $W = 24$ mois, pas de 1 mois,
produisant environ 171 fenêtres d'entraînement.

(5)~\textit{Calibration CP} : 20\% des observations ($\approx 34$ fenêtres)
réservées pour le calibration set, avec SSBC.

% ════════════════════════════════════════════════════════════════
\section{Conclusion}

\subsection{Bilan de la revue de littérature}

Ce document a présenté une analyse technique de treize travaux scientifiques
couvrant les fondations du VAE, les architectures récurrentes-variationnelles,
les Transformers spécialisés et les hybrides VAE-Transformer récents. L'analyse
est outillée par un tableau de benchmarks normalisés sur cinq datasets (SMD,
MSL, SMAP, SWaT, PSM) pour cinq modèles à résultats publics, et une matrice
qualitative comparative sur dix critères. Les benchmarks confirment que la
progression architecturale se traduit par des gains mesurables : OmniAnomaly
(F1 moyen 84.69\%) $\rightarrow$ Anomaly Transformer (94.91\%) $\rightarrow$
VAEAT (95.17\%) $\rightarrow$ MAAT (95.89\%).

\subsection{Les quatre contributions architecturales de BiVAT}

\textbf{D1  Encodeur hybride Transformer + BiLSTM.} Le Transformer seul est
sous-performant de 25 à 60\% sur $N < 300$ (Nakashima 2023) ; le BiLSTM seul
est supérieur sur séries financières courtes (MSE 0.408 vs 0.668) ; l'hybride
est systématiquement supérieur aux deux composantes isolées (VLT-Anomaly,
VAEAT).

\textbf{D2  Conformal Prediction.} Seule garantie formelle de coverage
distribution-free. MC Dropout exclu (overconfidence $N<200$) ; Deep Ensembles
exclus (coût $\times 5$ RAM). Adaptations SSBC, CP class-conditional, EnbPI
pour les contraintes BEAC.

\textbf{D3  SHAP inter-couches.} Stabilité 1.00 vs 0.20--0.81 LIME ; seule
méthode citée dans ICAAP/CCAR et rapports BIS FSI (2024) / ECB (2025) ;
précédent SHAPAttenVAE (2025).

\textbf{D4  Critère hybride reconstruction + AssDis.} Gain documenté de
$+18{,}82$ points F1 moyen (ablation Table~3, Xu et al.\ 2022) ; coût
computationnel nul car AssDis dérivée des attention maps existantes.

\subsection{Perspectives de déploiement à la BEAC}

BiVAT permettra de produire chaque mois un rapport structuré contenant :
(a)~une liste ordonnée des observations anormales avec scores de confiance et
intervalles CP ; (b)~une décomposition de la nature de chaque anomalie
(ponctuelle / contextuelle / collective) ; (c)~une attribution dimensionnelle
SHAP par indicateur IFS ; (d)~une recommandation de vérification calibrée selon
le niveau de confiance. Ce dispositif, conforme aux exigences COBAC et ancré
dans les recommandations BIS et FMI, constituerait le premier système de
surveillance macroprudentielle basé sur une architecture hybride VAE-Transformer
explicable et à incertitude formellement quantifiée déployé dans la zone CEMAC.

% ════════════════════════════════════════════════════════════════
\section*{Bibliographie}
\addcontentsline{toc}{section}{Bibliographie}

\begin{enumerate}

\item[\lbrack F1\rbrack] Kingma, D.P., \& Welling, M. (2013).
\textit{Auto-Encoding Variational Bayes}. arXiv:1312.6114.

\item[\lbrack F2\rbrack] An, J., \& Cho, S. (2015).
\textit{Variational Autoencoder Based Anomaly Detection Using Reconstruction
Probability}. SNU Data Mining Center Technical Report.

\item[\lbrack F3\rbrack] Su, Y., Zhao, Y., Niu, C., Liu, R., Sun, W., \&
Pei, D. (2019). Robust Anomaly Detection for Multivariate Time Series through
Stochastic Recurrent Neural Network. \textit{ACM SIGKDD}, pp.~2828--2837.
DOI:~10.1145/3292500.3330672.

\item[\lbrack T1\rbrack] Xu, J., Wu, H., Wang, J., \& Long, M. (2022).
Anomaly Transformer: Time Series Anomaly Detection with Association Discrepancy.
\textit{ICLR 2022 (Spotlight)}.
\url{https://ise.thss.tsinghua.edu.cn/~mlong/doc/anomaly-transformer-iclr22.pdf}

\item[\lbrack T2\rbrack] Tuli, S., Casale, G., \& Jennings, N.R. (2022).
TranAD: Deep Transformer Networks for Anomaly Detection in Multivariate Time
Series Data. \textit{PVLDB}, 15(6), 1201--1214. DOI:~10.14778/3514061.3514067.

\item[\lbrack H1\rbrack] Song, A., Seo, E., \& Kim, H. (2023). Anomaly
VAE-Transformer: A Deep Learning Approach for Anomaly Detection in
Decentralized Finance. \textit{IEEE Access}, 11, 98115--98131.
DOI:~10.1109/ACCESS.2023.3313448.

\item[\lbrack H2\rbrack] Li, C., Kiat, Y.C., Jing, J., \& Long, C. (2025).
T-VAE: Transformer-Based Variational AutoEncoder for Perceiving Anomalies in
Multivariate Time Series Data. \textit{Expert Systems}, 42, e70078.
DOI:~10.1111/exsy.70078.

\item[\lbrack H3\rbrack] Wang, H., Pi, D., Zhang, X., Liu, H., \& Guo, C.
(2022). Variational Transformer-based Anomaly Detection Approach for
Multivariate Time Series. \textit{Measurement}, 191, 110791.
DOI:~10.1016/j.measurement.2022.110791.

\item[\lbrack H4\rbrack] Zhang, H., Xia, Y., Yan, T., \& Liu, G. (2021).
Unsupervised Anomaly Detection in Multivariate Time Series through
Transformer-based Variational Autoencoder. \textit{CCDC 2021}, pp.~281--286.

\item[\lbrack A1\rbrack] He, X. et al. (2024). VAEAT: Variational AutoEncoder
with Adversarial Training for Multivariate Time Series Anomaly Detection.
\textit{Information Sciences}, 676, 120852. DOI:~10.1016/j.ins.2024.120852.

\item[\lbrack A2\rbrack] Xie, B., Huo, Z., \& Zhang, C. (2025). Unsupervised
Anomaly Detection in Time Series via Enhanced VAE-Transformer Framework.
\textit{CMC}, 84(1), 843--860. DOI:~10.32604/cmc.2025.063151.

\item[\lbrack A3\rbrack] Zhang, X., Wang, G., Chen, Y., Liu, Q., \& Wang, G.
(2025). Inter-layer Explainable Variational Autoencoder for Multivariate Time
Series Anomaly Detection (SHAPAttenVAE). \textit{Engineering Applications of
Artificial Intelligence}, 159, 111585. DOI:~10.1016/j.engappai.2025.111585.

\item[\lbrack A4\rbrack] Sellam, A.Z., Benaissa, I., Taleb-Ahmed, A.,
Patrono, L., \& Distante, C. (2025). Mamba Adaptive Anomaly Transformer (MAAT).
\textit{Engineering Applications of Artificial Intelligence}. arXiv:2502.07858.

\item[] Vaswani, A. et al. (2017). Attention Is All You Need.
\textit{NeurIPS}, 30, 5998--6008.

\item[] Pang, G., Shen, C., Cao, L., \& Van Den Hengel, A. (2021). Deep
Learning for Anomaly Detection: A Review. \textit{ACM Computing Surveys},
54(2), 1--38.

\item[] Vovk, V., Gammerman, A., \& Shafer, G. (2005).
\textit{Algorithmic Learning in a Random World}. Springer.

\item[] Gal, Y., \& Ghahramani, Z. (2016). Dropout as a Bayesian Approximation.
\textit{ICML}, pp.~1050--1059.

\item[] Lakshminarayanan, B., Pritzel, A., \& Blundell, C. (2017). Simple and
Scalable Predictive Uncertainty Estimation using Deep Ensembles. \textit{NeurIPS}, 30.

\item[] Lundberg, S.M., \& Lee, S.I. (2017). A Unified Approach to Interpreting
Model Predictions. \textit{NeurIPS}, 30.

\item[] Liu, Y., Wu, H., Wang, J., \& Long, M. (2022). Non-stationary
Transformers. \textit{NeurIPS}, 35. arXiv:2205.14415.

\item[] Chen, R., \& Xie, C. (2021). Conformal Prediction Intervals for Dynamic
Time-Series. arXiv:2010.09107.

\item[] Nakashima, K. et al. (2023). Benchmark Study on Anomaly Detection in
Time Series. \textit{PHM Society Annual Conference}.

\item[] BIS / IFC (2022). Machine learning applications in central banking.
\textit{IFC Bulletin~57}.
\url{https://www.bis.org/ifc/publ/ifcb57\_01\_rh.pdf}

\item[] BIS / Banque Nationale de Malaisie (2023). Machine learning for anomaly
detection in financial reporting. \textit{IFC Working Paper~23}.
\url{https://www.bis.org/ifc/publ/ifcwork23.pdf}

\item[] BIS (2024). Finding a needle in a haystack: ML for anomaly detection in
payment systems. \textit{BIS Working Paper~1188}.
\url{https://www.bis.org/ifc/publ/ifcb61\_26\_rh.pdf}

\item[] BIS (2024). Central banks and AI. \textit{BIS Annual Economic Report
2024, Chapter~3}.
\url{https://www.bis.org/publ/arpdf/ar2024e3.pdf}

\item[] BIS / FSI (2024). How regulators can address AI explainability.
\url{https://www.bis.org/fsi/fsipapers24.pdf}

\item[] ECB (2025). Guide to Internal Models (ML chapter).
\url{https://www.bankingsupervision.europa.eu}

\item[] PRA / Bank of England (2023). Model risk management principles for
banks. \textit{Supervisory Statement SS1/23}.
\url{https://www.bankofengland.co.uk/prudential-regulation/publication/2023}

\item[] IMF (2025). AI Projects in Financial Supervisory Authorities.
\textit{WP/25/199}.
\url{https://www.imf.org/-/media/files/publications/wp/2025/english/wpiea2025199-source-pdf.pdf}

\item[] Moukoko Ekambi, G.M. (2025). MKO-KAN : Détection de fraude Mobile Money
par réseaux de Kolmogorov-Arnold hybrides en zone CEMAC. Rapport de stage,
BEAC / ENSPD / Université de Douala, GISDIA Laboratory.

\end{enumerate}

\vspace{2cm}
\hrule
\vspace{0.5cm}
\noindent{\small\textit{Document rédigé dans le cadre du stage de recherche à
la Banque des États de l'Afrique Centrale (BEAC), M2 Data Science \& IA,
ENSPD / Université de Douala  GISDIA Laboratory. Août 2026.}}

\end{document}